% Paper 8, DRAFT 1 (home, 2026-11-10; owner 15:27 JST "ドラフトで書いておいて"; from notes/paper8/outline.md, the
% workplace skeleton of 2026-11-10). Not released. Numbers in the appendix tables come only from out/tables.tex
% (generated by notes/paper8/tables.py from committed records); numbers in the text are quoted from the committed
% records named beside them. Two device results are pending (marked PENDING): the N = 91 rerun with the ruler in
% every job (prereg/du91r_2026-11-10.md) and the Ising line against fold (prereg/isingL_2026-11-09.md).
\documentclass[10pt,a4paper]{article}
\usepackage[margin=2.2cm]{geometry}
\usepackage{amsmath,amssymb}
\usepackage{booktabs}
\usepackage{graphicx}
\usepackage[hidelinks]{hyperref}
\usepackage{url}
\setlength{\parskip}{4pt}
\setlength{\parindent}{0pt}
\sloppy
\newcommand{\pending}[1]{\textbf{[PENDING: #1]}}

\title{Wider and deeper without mitigation: IBM's dual-unitary chain at 91 qubits,\\ and three-dimensional lattices
through a placement-only OS}
\author{TN Holdings\\ \small contact: \texttt{love.os.architect@proton.me}}
\date{DRAFT 1, 10 November 2026 (not released)}

\begin{document}
\maketitle

\begin{abstract}
Paper 7 of this series carried user-written two-dimensional circuits through an operating layer (psf-os) that reads
no calibration-quality number, adds its own controls to each job and applies no error mitigation. Here the same
policy is carried to a chain that no fixed shape can hold and to three-dimensional lattices that the heavy-hex
coupling graph cannot embed. IBM's dual-unitary kicked-Ising chain (arXiv:2411.00765) at its largest width,
$N = 91$, read raw on ibm\_marrakesh: the decay rates agree with $\cos 2h$ at all three $h$ (largest
$|k - \cos 2h|$ 0.0037) and every depth up to $t = 24$ (2\,205 CZ) is readable, on a chain placed by a fixed rule
round the failed elements. \pending{second run, another device, ruler in every job.} In three dimensions, a place
rule that judges the place with the highest own mirror replaced a gate that had refused every place at 48 and 72
sites; kicked Ising on $4\times4\times3$ read at every angle and depth, Heisenberg XXX on $4\times2\times2$ read
with a N\'eel retention of 0.72 to 0.51. Two new layouts cut the transport CZ of the Heisenberg step by 42\%, and
were compared with the old one on the same places of two devices. On the way we changed our main measure for such
comparisons, from the retention of an order parameter to a normalised per-site error, because noise that slows the
dynamics can raise the retention. Pair creation in the Schwinger model read in 16 of 18 circuits and beat both
trivial guesses in all 18. A light-cone test failed as predicted: every front arrived late, while the shape of the
spread was reproduced. Every device result re-scores to the same numbers from its committed raw counts (21 of 21).
No quantum advantage is claimed: every answer is computed exactly on a classical machine.
\end{abstract}

\section{Introduction}\label{sec:intro}
Paper 7 (doi:10.5281/zenodo.24046673) reported user-written circuits on the two-dimensional lattice, a daily record
and IBM's dual-unitary chain at 51 and 71 sites, all under one policy: placement from the coupling graph and the
failed elements only, the device chosen by its queue, the job's own controls (floor, mirror, ring sentinel) in the
same submission, no mitigation, answers sealed by hash before each run. This paper keeps the policy unchanged and asks
whether it carries to
\begin{itemize}
\item a wider chain, where no fixed shape fits on the day's records and a fixed walk round the failed elements
(rule C) must place it;
\item three-dimensional lattices, which heavy-hex cannot embed, so that the OS must fold them and move qubits.
\end{itemize}
What is new in the OS: rule C for chains, a 3D layer with fold and line layouts, the place rule ``top'', and a
change of the main measure for layout comparisons (section~\ref{sec:layouts}). Not claimed: quantum advantage (every
answer is computed classically); that any mitigation works; that one layout is better in general (each comparison is
per place and day, with the measure fixed in advance).

\section{The dual-unitary chain at $N = 91$}\label{sec:du91}
\subsection{Placement}
On the records of 7 to 9 November no fixed chain shape of daily/du91 fits 91 sites on any of the three devices.
Rule C (\texttt{psf\_os.chain.place}, adopted on 6 November) walks round the failed elements in a fixed neighbour order
from fixed starts and takes the first fit; it is not a search, and the steps are recorded. On ibm\_marrakesh on 9
November it placed 91 sites in 1\,665 steps. Two checks preceded the run: the compiled circuits at $h = 0$ (Clifford)
were simulated at full size by the stabilizer method ($|\langle X\rangle| = 1$ on the light-cone site at every $t$),
and the closed form for $h > 0$ had been checked against statevectors and an MPS.

\subsection{Result (declared verdicts)}
\texttt{results/du91-2026-11-09}, prereg/du91\_2026-11-09.md. READS (MSE 0.0393 against a mean square 0.699). Decay
rates $k = 0.9956$, 0.9821, 0.9590 against $\cos 2h = 0.9950$, 0.9801, 0.9553: AGREES at all three $h$. Readable
depth 24 for every $h$. Sentinels: job A DEAD, job B CLEAN (as recorded). Table~1 sets the three widths side by side.

\subsection{After the fact}
The scatter of the rate fit is larger than shot noise ($\chi^2/5$ = 38, 26, 19 for the three $h$; with it the
deviations are 0.4 to 1.7 $\sigma$; \texttt{notes/du91\_2026-11-09}). In all three widths $h = 0.15$ came out slower
than $\cos 2h$, and in all three the ruler ($h = 0$) and $h = 0.15$ sat in different jobs.

\subsection{Second run}
\pending{prereg/du91r\_2026-11-10.md: fez or kingston, ruler $h = 0$ in both jobs (A: 0, 0.05, 0.10; B: 0, 0.15),
each rate against the ruler of its own job; the cross-job rate and the ruler ratio beside.}

\section{The three-dimensional layer of the OS}\label{sec:os3d}
\subsection{Layouts}
A lattice $L\times M\times N$ is folded onto heavy-hex (fold v3.1): blocks are placed by translation and the bonds
between them are carried by swaps over bridge qubits, on fixed schedules found offline and kept as constants. Two
layouts followed: \emph{carry} (an $x$ bond waits for a later swap of the same step and rides on it) and \emph{line}
(two blocks interleaved on one line, no empty swap, even steps reversed, the $y$ bonds and line operations at step
boundaries merged). Table~4 gives the CZ per step.

\subsection{Choosing the place}
A pre-run gate (probes for floor marks and low mirror sites) worked at 16 sites: on 7 November on ibm\_marrakesh it
passed one of four places, which was judged (\texttt{results/os3d-gate-2026-11-07}). At 48 sites it refused every one
of five places on ibm\_fez (all had $|1\rangle$ floor marks), so no 3D arm was judged
(\texttt{results/os3d-mid-2026-11-07}). The owner replaced the stop by the rule ``top'': the place whose own $t = 1$
mirror is highest is judged, and the gate's marks are recorded beside. On ibm\_kingston the judged $4\times4\times3$
place read at $t = 1$--3 (retention 0.92, 0.95, 0.95) and had the smallest error of the five places; the gate alone
would have refused it and every other place. $4\times6\times3$ (72 sites, one place) read with retention 0.88, 0.88,
0.86 (\texttt{results/os3d-kingston-top-2026-11-07}).

\section{Physics A: kicked Ising on 3D lattices}\label{sec:physA}
\texttt{results/os3d-physA-2026-11-07}, ibm\_fez, rule top, $\theta = 0.4$, 0.6, 0.8. $4\times4\times3$ (48 sites,
CZ up to 1\,464) READS at every $\theta$ and $t = 1$--3, with retention 0.73 to 0.83; $4\times2\times2$ READS at
every $\theta$ and $t = 1$--6 (Table~2). The magnetisation DIFFERS on most rows under the absolute line of 0.05; as
in paper 7, the retention beside it shows how much signal was kept.

\section{Physics B: Heisenberg XXX on $4\times2\times2$}\label{sec:physB}
\texttt{results/os3d-heisB-2026-11-08}, ibm\_fez. READS at $t = 1$--4, DIFFERS; N\'eel retention 0.72, 0.62, 0.51,
0.53 at 176 to 692 CZ (Table~3).

\section{Layouts compared on the same places}\label{sec:layouts}
\subsection{Transport CZ}
Old, carry and line give 176, 152 and 116 CZ at $t = 1$ and 692, 596 and 404 at $t = 4$ for the Heisenberg step
(42\% fewer at $t = 4$ with line); for Ising, line gives 60 against fold's 108 per step at $4\times2\times2$ and 126
against 234 at $4\times3\times2$ (Table~4). Every layout was checked against the exact answer before any run.

\subsection{Declared comparison}
The three Heisenberg layouts ran in one job at the same places on ibm\_marrakesh and on ibm\_kingston (8 November;
Table~5). The declared measure was the difference of N\'eel retention, line minus old, with an expectation of $\ge 0$.
At the judged place on marrakesh: $+0.039$, $+0.079$, $+0.096$, $+0.077$ at $t = 1$--4 (resolved at $t = 2$ and 3).
On kingston: $+0.022$, $+0.056$, $+0.050$, $-0.051$; the expectation failed at $t = 4$.

\subsection{Why the measure changed}
A retention can rise for the wrong reason: noise that slows the exchange keeps the N\'eel order longer. From 8
November the main measure for layout comparisons is the normalised per-site error
$\mathrm{nMSE} = \overline{(z - e)^2}/\overline{e^2}$, with the retention beside it
(\texttt{notes/reporting/layout\_comparison\_metric.md}). Applied after the fact to the same arms, line has the lower
nMSE in 28 of 32 place-and-$t$ cells (Table~5). The two runs keep their declared verdicts.

\subsection{Ising: line against fold}
\pending{prereg/isingL\_2026-11-09.md; the declared noiseless pre-run check passed on 10 November: all 48 arms READ and
AGREE, \texttt{notes/isingL\_noiseless\_2026-11-10}.}

\section{A new problem: pair creation in the Schwinger model}\label{sec:schwinger}
\texttt{results/schw-2026-11-09}, ibm\_marrakesh; Martinez et al.\ (2016) form, $N = 6$ and 8, parameter sets M1--M3,
1, 2 and 5 Trotter steps, a swap network on a chain. READS 16 of 18, BEATS (both the bare vacuum and the all-zero
guess) 18 of 18, particle density AGREES within 0.05 in 14 of 18. The two that did not read are $N = 8$ at 5 steps
for M1 and M3 (420 CZ). The total charge $\sum Z$ (exactly 0) is recorded raw, never used to correct; it drifted to
$+0.1$ to $+0.6$ (Table~6). An estimate after the run (\texttt{notes/schwinger\_scale\_2026-11-10}): the swap network
costs $3(N-1)$ CZ per qubit per step, so larger $N$ would not read with this circuit.

\section{The light cone in 3D, retried on a device}\label{sec:lc}
\texttt{results/lc3d-2026-11-09}, ibm\_marrakesh: one flipped spin in the all-up $4\times2\times2$ XXX lattice, with
a no-flip arm in the same job to subtract the background. The test had been dropped on an estimate that noise flips
would bury the signal; it was run after all, with the failure beyond distance 1 predicted. Every front DIFFERS at the
judged place: distance 1 arrived at $t = 2$ (exact $t = 1$), distance 2 at $t = 6$ (exact 3), distance 3 not within
$t = 6$. The background held 2.1 to 3.9 flipped spins and the excitation kept 0.76 down to 0.43 of itself, while the
shape of the spread matched the exact one within a total variation distance of 0.025 to 0.144 (Table~8).

\section{The simulator and the device}\label{sec:sim}
A trajectory simulator with that day's calibrated noise, never fitted to the device
(\texttt{notes/os3d\_sim\_2026-11-07/realsim.py}), was compared once with the device
(\texttt{notes/simproto\_2026-11-08}). For kicked Ising $4\times2\times2$ the simulated value over the device value
had a mean of 0.987 (0.929 to 1.025); for Heisenberg it overestimated the retention by 15\% to 28\% (Table~7). The
difference between layouts, not the level, is what the comparisons above need.

\section{Limits and lessons}\label{sec:limits}
\begin{itemize}
\item \emph{Two mirrors} (a second, twirled mirror) were not added: a planted error of 0.02 could not be told apart
within 1.5 times the device time; a hypothesis about coherent CZ errors was withdrawn after a computing error was
found (\texttt{notes/two\_mirrors\_2026-11-08}).
\item \emph{Three mirrors} usable at $t = 1$ only; at $t \ge 2$ the product rule they rest on fails under gate-local
noise alone (\texttt{notes/three\_mirrors\_2026-11-09}).
\item The retention bias that led to nMSE (section~\ref{sec:layouts}); a declared expectation that failed (kingston,
$t = 4$); the fit scatter at $N = 91$ and the slow $h = 0.15$ (section~\ref{sec:du91}).
\item Three devices of one family and a few days; the 3D lattices are small ($\le 72$ sites) and every answer is
classically computable.
\end{itemize}

\section{Reproduction}\label{sec:repro}
\texttt{python reproduce.py} re-runs every committed scorer, unchanged, on a copy of each result from its committed
raw counts and compares every number with the committed scores. For the 21 device results of papers 7 and 8 it prints
SAME for all 21; for three Heisenberg results whose scorer gained keys later, the scorer as committed at the time is
taken from the history. Answers files are regenerated by committed code and their hashes compared with those
committed before each run; job folders and job identifiers are committed.

\section{Conclusion}
The no-map, no-mitigation policy carried to IBM's 91-site dual-unitary chain, placed by a fixed walk round failed
elements, with decay rates within 0.004 of $\cos 2h$; and to 3D lattices up to 72 sites through a fold of heavy-hex,
with the place chosen by its own mirror. Layouts that cut the transport CZ by 42\% were compared on the same places of
two devices, and the comparison taught us to measure per-site error rather than retention. A new problem read, and a
light-cone test failed as predicted. No advantage is claimed. \pending{second $N = 91$ run; Ising line.}

\section*{Acknowledgements}
We acknowledge the use of IBM Quantum services for this work. The views expressed are those of the authors and do not
reflect the official policy or position of IBM or the IBM Quantum team.

\section*{Data and code}
TOQB (\url{https://github.com/TN-Holdings-LLC/TOQB}) holds the declarations (\texttt{prereg/}), the results with raw
counts, job identifiers and released answers (\texttt{results/}), the scorers (\texttt{gpu/}), \texttt{reproduce.py}
and the tables of this paper (\texttt{notes/paper8/tables.py}). psf-os (\url{https://github.com/TN-Holdings-LLC/psf-os})
is the operating layer. \pending{commit hashes at release.}

\section*{References}
\begin{enumerate}\small
\item IBM Quantum, dual-unitary kicked-Ising experiment at 51, 71 and 91 qubits, arXiv:2411.00765.
\item E.~A.~Martinez et al., Real-time dynamics of lattice gauge theories with a few-qubit quantum computer, Nature
534, 516 (2016), arXiv:1605.04570.
\item TN Holdings, paper 7 of this series, doi:10.5281/zenodo.24046673.
\item TN Holdings, paper 6 of this series, doi:10.5281/zenodo.23883938.
\end{enumerate}

\appendix
\section{Generated tables}
Generated by \texttt{notes/paper8/tables.py} from the committed records; no number in them is typed by hand.
{\setlength{\tabcolsep}{3pt}
% generated by notes/paper8/tables.py from committed records (DRAFT); do not edit by hand
\begin{table}[htbp]\centering\small
\caption{The dual-unitary chain (arXiv 2411.00765) raw: N = 51, 71 (paper 7) and 91 twice (results/du51-2026-11-07, du71-2026-11-07, du91-2026-11-09, night-2026-11-10/du91r; in the second run each rate is against the ruler of its own job). From the committed records named in each caption; raw device values, no mitigation; verdicts as declared before each run; after-the-fact numbers are marked so.}
\resizebox{\ifdim\width>\textwidth\textwidth\else\width\fi}{!}{\begin{tabular}{lllllllll}\toprule
N & device & placement & CZ (t 24) & READS (MSE / mean sq.) & k - cos2h, h 0.05 & h 0.10 & h 0.15 & readable depth (h 0 / .05 / .10 / .15) \\ \midrule
51 & ibm\_kingston & shape S1 & 1225 & READS (0.0403 / 0.699) & +0.0000 AGREES & +0.0020 AGREES & +0.0045 AGREES & 24 / 24 / 24 / 24 \\
71 & ibm\_kingston & shape S1 & 1715 & READS (0.0508 / 0.699) & -0.0002 AGREES & -0.0002 AGREES & +0.0017 AGREES & 24 / 24 / 24 / 24 \\
91 & ibm\_marrakesh & C & 2205 & READS (0.0393 / 0.699) & +0.0006 AGREES & +0.0020 AGREES & +0.0037 AGREES & 24 / 24 / 24 / 24 \\
91 & ibm\_fez (2nd run, ruler in both jobs) & C & 2205 & READS (0.2872 / 0.760) & +0.0010 AGREES & +0.0013 AGREES & +0.0037 AGREES & 24 / 24 / 16 / 24 \\
\bottomrule\end{tabular}}
\end{table}

\begin{table}[htbp]\centering\small
\caption{Physics A: kicked Ising on 3D lattices, ibm\_fez 2026-11-07, the judged place (rule top), 3D sites (results/os3d-physA-2026-11-07). From the committed records named in each caption; raw device values, no mitigation; verdicts as declared before each run; after-the-fact numbers are marked so.}
\resizebox{\ifdim\width>\textwidth\textwidth\else\width\fi}{!}{\begin{tabular}{llllllllll}\toprule
lattice & theta & t & CZ & M & M exact & retention & READS & AGREES & BEATS \\ \midrule
4x4x3 & 0.4 & 1 & 488 & 0.729 & 0.921 & 0.791 & READS & DIFFERS & DOES NOT BEAT \\
4x4x3 & 0.4 & 2 & 976 & 0.638 & 0.811 & 0.787 & READS & DIFFERS & DOES NOT BEAT \\
4x4x3 & 0.4 & 3 & 1464 & 0.534 & 0.729 & 0.733 & READS & DIFFERS & BEATS \\
4x4x3 & 0.6 & 1 & 488 & 0.659 & 0.825 & 0.798 & READS & DIFFERS & DOES NOT BEAT \\
4x4x3 & 0.6 & 2 & 976 & 0.503 & 0.628 & 0.800 & READS & DIFFERS & BEATS \\
4x4x3 & 0.6 & 3 & 1464 & 0.393 & 0.512 & 0.767 & READS & DIFFERS & BEATS \\
4x4x3 & 0.8 & 1 & 488 & 0.556 & 0.697 & 0.798 & READS & DIFFERS & BEATS \\
4x4x3 & 0.8 & 2 & 976 & 0.367 & 0.440 & 0.834 & READS & DIFFERS & BEATS \\
4x4x3 & 0.8 & 3 & 1464 & 0.255 & 0.312 & 0.818 & READS & DIFFERS & BEATS \\
4x2x2 & 0.4 & 1 & 108 & 0.829 & 0.921 & 0.900 & READS & DIFFERS & DOES NOT BEAT \\
4x2x2 & 0.4 & 2 & 216 & 0.700 & 0.794 & 0.882 & READS & DIFFERS & BEATS \\
4x2x2 & 0.4 & 3 & 324 & 0.616 & 0.717 & 0.859 & READS & DIFFERS & BEATS \\
4x2x2 & 0.4 & 4 & 432 & 0.548 & 0.634 & 0.863 & READS & DIFFERS & BEATS \\
4x2x2 & 0.4 & 5 & 540 & 0.463 & 0.537 & 0.862 & READS & DIFFERS & BEATS \\
4x2x2 & 0.4 & 6 & 648 & 0.408 & 0.492 & 0.830 & READS & DIFFERS & BEATS \\
4x2x2 & 0.6 & 1 & 108 & 0.744 & 0.825 & 0.901 & READS & DIFFERS & BEATS \\
4x2x2 & 0.6 & 2 & 216 & 0.544 & 0.607 & 0.895 & READS & DIFFERS & BEATS \\
4x2x2 & 0.6 & 3 & 324 & 0.450 & 0.513 & 0.877 & READS & DIFFERS & BEATS \\
4x2x2 & 0.6 & 4 & 432 & 0.380 & 0.448 & 0.848 & READS & DIFFERS & BEATS \\
4x2x2 & 0.6 & 5 & 540 & 0.326 & 0.418 & 0.779 & READS & DIFFERS & BEATS \\
4x2x2 & 0.6 & 6 & 648 & 0.284 & 0.429 & 0.661 & READS & DIFFERS & BEATS \\
4x2x2 & 0.8 & 1 & 108 & 0.630 & 0.697 & 0.904 & READS & DIFFERS & BEATS \\
4x2x2 & 0.8 & 2 & 216 & 0.390 & 0.425 & 0.918 & READS & AGREES & BEATS \\
4x2x2 & 0.8 & 3 & 324 & 0.299 & 0.321 & 0.932 & READS & AGREES & BEATS \\
4x2x2 & 0.8 & 4 & 432 & 0.232 & 0.304 & 0.764 & READS & DIFFERS & BEATS \\
4x2x2 & 0.8 & 5 & 540 & 0.188 & 0.301 & 0.625 & READS & DIFFERS & BEATS \\
4x2x2 & 0.8 & 6 & 648 & 0.152 & 0.278 & 0.547 & READS & DIFFERS & BEATS \\
\bottomrule\end{tabular}}
\end{table}

\begin{table}[htbp]\centering\small
\caption{Physics B: Heisenberg XXX 4x2x2, ibm\_fez 2026-11-08, the judged place 0 (results/os3d-heisB-2026-11-08/fez). From the committed records named in each caption; raw device values, no mitigation; verdicts as declared before each run; after-the-fact numbers are marked so.}
\resizebox{\ifdim\width>\textwidth\textwidth\else\width\fi}{!}{\begin{tabular}{llllllll}\toprule
t & CZ & m\_s & m\_s exact & retention & nMSE (after the fact) & READS & AGREES \\ \midrule
1 & 176 & 0.669 & 0.932 & 0.719 & 0.088 & READS & DIFFERS \\
2 & 348 & 0.469 & 0.757 & 0.619 & 0.163 & READS & DIFFERS \\
3 & 520 & 0.270 & 0.530 & 0.509 & 0.299 & READS & DIFFERS \\
4 & 692 & 0.164 & 0.308 & 0.532 & 0.369 & READS & DIFFERS \\
\bottomrule\end{tabular}}
\end{table}

\begin{table}[htbp]\centering\small
\caption{CZ per layout and t (Heisenberg: heisE marrakesh build, judged place; Ising: the isingL build of the noiseless check, place 0; notes/isingL\_noiseless\_2026-11-10). From the committed records named in each caption; raw device values, no mitigation; verdicts as declared before each run; after-the-fact numbers are marked so.}
\resizebox{\ifdim\width>\textwidth\textwidth\else\width\fi}{!}{\begin{tabular}{llllll}\toprule
model & layout & CZ t 1 & t 2 & t 3 & t 4 \\ \midrule
Heisenberg 4x2x2 & old & 176 & 348 & 520 & 692 \\
Heisenberg 4x2x2 & carry & 152 & 300 & 448 & 596 \\
Heisenberg 4x2x2 & line & 116 & 212 & 308 & 404 \\
Ising 4x2x2 & fold & 108 & 216 & 324 & 432 \\
Ising 4x2x2 & line & 60 & 120 & 180 & 240 \\
Ising 4x3x2 & fold & 234 & 468 & 702 & 936 \\
Ising 4x3x2 & line & 126 & 252 & 378 & 504 \\
\bottomrule\end{tabular}}
\end{table}

\begin{table}[htbp]\centering\small
\caption{Heisenberg layouts on the same places: line against old (results/os3d-heisE-2026-11-08, os3d-heisE-kingston-2026-11-08); the declared measure was the retention difference, nMSE is after the fact (notes/reporting/layout\_comparison\_metric.md). From the committed records named in each caption; raw device values, no mitigation; verdicts as declared before each run; after-the-fact numbers are marked so.}
\resizebox{\ifdim\width>\textwidth\textwidth\else\width\fi}{!}{\begin{tabular}{llllllll}\toprule
device & place & t & CZ old / line & retention line - old (declared) & resolved & nMSE old & nMSE line (after the fact) \\ \midrule
ibm\_marrakesh & 0 & 1 & 176 / 116 & +0.038 & no & 0.083 & 0.064 \\
ibm\_marrakesh & 0 & 2 & 348 / 212 & +0.062 & yes & 0.154 & 0.100 \\
ibm\_marrakesh & 0 & 3 & 520 / 308 & -0.016 & no & 0.261 & 0.270 \\
ibm\_marrakesh & 0 & 4 & 692 / 404 & +0.009 & no & 0.449 & 0.322 \\
ibm\_marrakesh & 1 (judged) & 1 & 176 / 116 & +0.039 & no & 0.046 & 0.031 \\
ibm\_marrakesh & 1 (judged) & 2 & 348 / 212 & +0.079 & yes & 0.118 & 0.075 \\
ibm\_marrakesh & 1 (judged) & 3 & 520 / 308 & +0.096 & yes & 0.229 & 0.132 \\
ibm\_marrakesh & 1 (judged) & 4 & 692 / 404 & +0.077 & no & 0.425 & 0.267 \\
ibm\_marrakesh & 2 & 1 & 176 / 116 & +0.037 & no & 0.055 & 0.042 \\
ibm\_marrakesh & 2 & 2 & 348 / 212 & +0.057 & no & 0.124 & 0.091 \\
ibm\_marrakesh & 2 & 3 & 520 / 308 & +0.094 & yes & 0.228 & 0.152 \\
ibm\_marrakesh & 2 & 4 & 692 / 404 & +0.111 & no & 0.321 & 0.236 \\
ibm\_marrakesh & 3 & 1 & 212 / 128 & +0.012 & no & 0.138 & 0.128 \\
ibm\_marrakesh & 3 & 2 & 420 / 236 & +0.070 & yes & 0.193 & 0.130 \\
ibm\_marrakesh & 3 & 3 & 628 / 344 & -0.053 & no & 0.377 & 0.437 \\
ibm\_marrakesh & 3 & 4 & 836 / 452 & -0.037 & no & 0.344 & 0.518 \\
ibm\_kingston & 0 & 1 & 176 / 116 & +0.029 & no & 0.046 & 0.034 \\
ibm\_kingston & 0 & 2 & 348 / 212 & +0.044 & no & 0.118 & 0.091 \\
ibm\_kingston & 0 & 3 & 520 / 308 & +0.000 & no & 0.223 & 0.212 \\
ibm\_kingston & 0 & 4 & 692 / 404 & -0.000 & no & 0.429 & 0.409 \\
ibm\_kingston & 1 (judged) & 1 & 176 / 116 & +0.022 & no & 0.043 & 0.035 \\
ibm\_kingston & 1 (judged) & 2 & 348 / 212 & +0.056 & no & 0.128 & 0.090 \\
ibm\_kingston & 1 (judged) & 3 & 520 / 308 & +0.050 & no & 0.250 & 0.186 \\
ibm\_kingston & 1 (judged) & 4 & 692 / 404 & -0.051 & no & 0.466 & 0.455 \\
ibm\_kingston & 2 & 1 & 176 / 116 & +0.057 & yes & 0.049 & 0.026 \\
ibm\_kingston & 2 & 2 & 348 / 212 & +0.109 & yes & 0.142 & 0.070 \\
ibm\_kingston & 2 & 3 & 520 / 308 & +0.144 & yes & 0.302 & 0.151 \\
ibm\_kingston & 2 & 4 & 692 / 404 & +0.109 & no & 0.471 & 0.394 \\
ibm\_kingston & 3 & 1 & 176 / 116 & +0.080 & yes & 0.135 & 0.056 \\
ibm\_kingston & 3 & 2 & 348 / 212 & +0.057 & no & 0.233 & 0.183 \\
ibm\_kingston & 3 & 3 & 520 / 308 & +0.096 & yes & 0.366 & 0.232 \\
ibm\_kingston & 3 & 4 & 692 / 404 & +0.001 & no & 0.388 & 0.444 \\
\bottomrule\end{tabular}}
\end{table}

\begin{table}[htbp]\centering\small
\caption{Schwinger pair creation from the bare vacuum, ibm\_marrakesh 2026-11-09, sentinel CLEAN (results/schw-2026-11-09). From the committed records named in each caption; raw device values, no mitigation; verdicts as declared before each run; after-the-fact numbers are marked so.}
\resizebox{\ifdim\width>\textwidth\textwidth\else\width\fi}{!}{\begin{tabular}{llllllllll}\toprule
N & set & steps & CZ & nMSE & READS & BEATS & nu / exact & nu & Sum Z (exact 0) \\ \midrule
6 & M1 & 1 & 45 & 0.009 & READS & BEATS & 0.133 / 0.097 & AGREES & +0.124 \\
6 & M1 & 2 & 90 & 0.029 & READS & BEATS & 0.332 / 0.315 & AGREES & +0.230 \\
6 & M1 & 5 & 225 & 0.446 & READS & BEATS & 0.443 / 0.333 & DIFFERS & +0.406 \\
6 & M2 & 1 & 45 & 0.013 & READS & BEATS & 0.140 / 0.097 & AGREES & +0.154 \\
6 & M2 & 2 & 90 & 0.020 & READS & BEATS & 0.331 / 0.327 & AGREES & +0.214 \\
6 & M2 & 5 & 225 & 0.188 & READS & BEATS & 0.574 / 0.563 & AGREES & +0.418 \\
6 & M3 & 1 & 45 & 0.008 & READS & BEATS & 0.130 / 0.097 & AGREES & +0.115 \\
6 & M3 & 2 & 90 & 0.022 & READS & BEATS & 0.333 / 0.318 & AGREES & +0.215 \\
6 & M3 & 5 & 225 & 0.384 & READS & BEATS & 0.446 / 0.361 & DIFFERS & +0.367 \\
8 & M1 & 1 & 84 & 0.015 & READS & BEATS & 0.146 / 0.101 & AGREES & +0.206 \\
8 & M1 & 2 & 168 & 0.036 & READS & BEATS & 0.345 / 0.328 & AGREES & +0.370 \\
8 & M1 & 5 & 420 & 0.677 & DOES NOT READ & BEATS & 0.467 / 0.328 & DIFFERS & +0.612 \\
8 & M2 & 1 & 84 & 0.013 & READS & BEATS & 0.144 / 0.101 & AGREES & +0.186 \\
8 & M2 & 2 & 168 & 0.046 & READS & BEATS & 0.343 / 0.339 & AGREES & +0.369 \\
8 & M2 & 5 & 420 & 0.234 & READS & BEATS & 0.569 / 0.557 & AGREES & +0.532 \\
8 & M3 & 1 & 84 & 0.016 & READS & BEATS & 0.148 / 0.101 & AGREES & +0.259 \\
8 & M3 & 2 & 168 & 0.044 & READS & BEATS & 0.344 / 0.331 & AGREES & +0.366 \\
8 & M3 & 5 & 420 & 0.566 & DOES NOT READ & BEATS & 0.460 / 0.359 & DIFFERS & +0.539 \\
\bottomrule\end{tabular}}
\end{table}

\begin{table}[htbp]\centering\small
\caption{The simulator (notes/os3d\_sim\_2026-11-07/realsim.py --expect; noise from the device's calibration of that day, never fitted) against the device (notes/simproto\_2026-11-08). From the committed records named in each caption; raw device values, no mitigation; verdicts as declared before each run; after-the-fact numbers are marked so.}
\resizebox{\ifdim\width>\textwidth\textwidth\else\width\fi}{!}{\begin{tabular}{lllllll}\toprule
model & arm & CZ & exact & device & simulator (+- SE) & simulator / device retention \\ \midrule
Heisenberg B & h3d0\_t1 & 176 & 0.932 & 0.669 & 0.771 +- 0.009 & 1.152 \\
Heisenberg B & h3d0\_t2 & 348 & 0.757 & 0.469 & 0.551 +- 0.009 & 1.176 \\
Heisenberg B & h3d0\_t3 & 520 & 0.530 & 0.270 & 0.345 +- 0.009 & 1.278 \\
Heisenberg B & h3d0\_t4 & 692 & 0.308 & 0.164 & 0.196 +- 0.008 & 1.198 \\
Ising A 4x2x2 & g3d3\_th04\_t1 & 108 & 0.921 & 0.829 & 0.849 +- 0.006 & 1.025 \\
Ising A 4x2x2 & g3d3\_th04\_t2 & 216 & 0.794 & 0.700 & 0.704 +- 0.007 & 1.006 \\
Ising A 4x2x2 & g3d3\_th04\_t3 & 324 & 0.717 & 0.616 & 0.622 +- 0.007 & 1.010 \\
Ising A 4x2x2 & g3d3\_th04\_t4 & 432 & 0.634 & 0.548 & 0.539 +- 0.009 & 0.984 \\
Ising A 4x2x2 & g3d3\_th04\_t5 & 540 & 0.537 & 0.463 & 0.463 +- 0.008 & 1.002 \\
Ising A 4x2x2 & g3d3\_th04\_t6 & 648 & 0.492 & 0.408 & 0.408 +- 0.008 & 1.000 \\
Ising A 4x2x2 & g3d3\_th06\_t1 & 108 & 0.825 & 0.744 & 0.755 +- 0.005 & 1.015 \\
Ising A 4x2x2 & g3d3\_th06\_t2 & 216 & 0.607 & 0.544 & 0.555 +- 0.005 & 1.021 \\
Ising A 4x2x2 & g3d3\_th06\_t3 & 324 & 0.513 & 0.450 & 0.454 +- 0.006 & 1.009 \\
Ising A 4x2x2 & g3d3\_th06\_t4 & 432 & 0.448 & 0.380 & 0.384 +- 0.006 & 1.009 \\
Ising A 4x2x2 & g3d3\_th06\_t5 & 540 & 0.418 & 0.326 & 0.306 +- 0.006 & 0.939 \\
Ising A 4x2x2 & g3d3\_th06\_t6 & 648 & 0.429 & 0.284 & 0.264 +- 0.006 & 0.929 \\
Ising A 4x2x2 & g3d3\_th08\_t1 & 108 & 0.697 & 0.630 & 0.637 +- 0.005 & 1.012 \\
Ising A 4x2x2 & g3d3\_th08\_t2 & 216 & 0.425 & 0.390 & 0.389 +- 0.003 & 0.999 \\
Ising A 4x2x2 & g3d3\_th08\_t3 & 324 & 0.321 & 0.299 & 0.286 +- 0.004 & 0.957 \\
Ising A 4x2x2 & g3d3\_th08\_t4 & 432 & 0.304 & 0.232 & 0.225 +- 0.005 & 0.968 \\
Ising A 4x2x2 & g3d3\_th08\_t5 & 540 & 0.301 & 0.188 & 0.179 +- 0.005 & 0.952 \\
Ising A 4x2x2 & g3d3\_th08\_t6 & 648 & 0.278 & 0.152 & 0.141 +- 0.004 & 0.930 \\
\bottomrule\end{tabular}}
\end{table}

\begin{table}[htbp]\centering\small
\caption{The light cone of one flipped spin, 4x2x2 XXX, ibm\_marrakesh 2026-11-09, judged place 0; fronts: d 1: exact t 1, measured 2 (FRONT DIFFERS); d 2: exact t 3, measured 6 (FRONT DIFFERS); d 3: exact t 6, measured none (FRONT DIFFERS) (results/lc3d-2026-11-09). From the committed records named in each caption; raw device values, no mitigation; verdicts as declared before each run; after-the-fact numbers are marked so.}
\resizebox{\ifdim\width>\textwidth\textwidth\else\width\fi}{!}{\begin{tabular}{lllll}\toprule
t & CZ & excitation kept K & background flips & shape TVD \\ \midrule
1 & 116 & 0.758 & 2.317 & 0.025 \\
2 & 212 & 0.646 & 2.147 & 0.056 \\
3 & 308 & 0.614 & 2.879 & 0.049 \\
4 & 404 & 0.572 & 3.099 & 0.053 \\
5 & 500 & 0.441 & 3.621 & 0.144 \\
6 & 596 & 0.432 & 3.876 & 0.102 \\
\bottomrule\end{tabular}}
\end{table}

\begin{table}[htbp]\centering\small
\caption{Ising (theta 0.6): the line layout against the fold at the same places, ibm\_fez 2026-11-10; main measure nMSE as declared (results/night-2026-11-10/isingL). From the committed records named in each caption; raw device values, no mitigation; verdicts as declared before each run; after-the-fact numbers are marked so.}
\resizebox{\ifdim\width>\textwidth\textwidth\else\width\fi}{!}{\begin{tabular}{llllllll}\toprule
lattice & place & t & CZ fold / line & nMSE fold & nMSE line & line - fold & resolved \\ \midrule
4x2x2 & 0 & 1 & 108 / 60 & 0.0149 & 0.0141 & -0.0008 & no \\
4x2x2 & 0 & 2 & 216 / 120 & 0.0239 & 0.0300 & +0.0060 & yes \\
4x2x2 & 0 & 3 & 324 / 180 & 0.0647 & 0.0944 & +0.0297 & yes \\
4x2x2 & 0 & 4 & 432 / 240 & 0.1045 & 0.1498 & +0.0454 & yes \\
4x2x2 & 1 & 1 & 108 / 60 & 0.0191 & 0.0118 & -0.0074 & yes \\
4x2x2 & 1 & 2 & 216 / 120 & 0.0227 & 0.0175 & -0.0051 & yes \\
4x2x2 & 1 & 3 & 324 / 180 & 0.0954 & 0.0381 & -0.0574 & yes \\
4x2x2 & 1 & 4 & 432 / 240 & 0.1852 & 0.0776 & -0.1076 & yes \\
4x2x2 & 2 & 1 & 108 / 60 & 0.0169 & 0.0191 & +0.0022 & no \\
4x2x2 & 2 & 2 & 216 / 120 & 0.0295 & 0.0287 & -0.0008 & no \\
4x2x2 & 2 & 3 & 324 / 180 & 0.0581 & 0.0476 & -0.0105 & yes \\
4x2x2 & 2 & 4 & 432 / 240 & 0.1001 & 0.0609 & -0.0392 & yes \\
4x2x2 & 3 (judged) & 1 & 108 / 60 & 0.0145 & 0.0139 & -0.0007 & no \\
4x2x2 & 3 (judged) & 2 & 216 / 120 & 0.0189 & 0.0166 & -0.0022 & no \\
4x2x2 & 3 (judged) & 3 & 324 / 180 & 0.0406 & 0.0323 & -0.0082 & yes \\
4x2x2 & 3 (judged) & 4 & 432 / 240 & 0.0732 & 0.0487 & -0.0244 & yes \\
4x3x2 & 0 (judged) & 1 & 234 / 126 & 0.0332 & 0.0358 & +0.0026 & no \\
4x3x2 & 0 (judged) & 2 & 468 / 252 & 0.0476 & 0.1248 & +0.0772 & yes \\
4x3x2 & 0 (judged) & 3 & 702 / 378 & 0.1072 & 0.1884 & +0.0812 & yes \\
4x3x2 & 0 (judged) & 4 & 936 / 504 & 0.2038 & 0.2477 & +0.0439 & yes \\
4x3x2 & 1 & 1 & 234 / 126 & 0.0334 & 0.0644 & +0.0310 & yes \\
4x3x2 & 1 & 2 & 468 / 252 & 0.0433 & 0.1173 & +0.0740 & yes \\
4x3x2 & 1 & 3 & 702 / 378 & 0.1349 & 0.2832 & +0.1482 & yes \\
4x3x2 & 1 & 4 & 936 / 504 & 0.2696 & 0.3897 & +0.1201 & yes \\
\bottomrule\end{tabular}}
\end{table}

\begin{table}[htbp]\centering\small
\caption{IBM's 2023 utility observables (Kim et al. 2023) on the Eagle graph embedded in ibm\_fez, 2026-11-10, raw, 4000 shots, 17 values of theta\_h; sentinel DEAD (results/utility-2026-11-10). From the committed records named in each caption; raw device values, no mitigation; verdicts as declared before each run; after-the-fact numbers are marked so.}
\resizebox{\ifdim\width>\textwidth\textwidth\else\width\fi}{!}{\begin{tabular}{lllllll}\toprule
observable & embedding & CZ & verdict (MSE / mean sq.) & AGREES of 17 & theta\_h = pi/2: raw / exact & floor / mirror parity \\ \midrule
Z\_106, 4 steps (Fig. 2) & 0 & 564 & READS (0.0207 / 0.4145) & 7 & +0.14 / +0.00 & 0.991 / 0.504 \\
weight-10, 5 steps (Fig. 3) & 0 & 705 & READS (0.0226 / 0.1672) & 11 & +0.61 / +1.00 & 0.852 / 0.122 \\
weight-17, 5 steps (Fig. 3) & 1 & 690 & DOES NOT READ (0.1091 / 0.1348) & 12 & -0.09 / -1.00 & 0.310 / -0.004 \\
Fig. 4a product, 5 steps + R\_X & 1 & 690 & DOES NOT READ (0.1023 / 0.1200) & 12 & -0.07 / -1.00 & 0.305 / -0.009 \\
\bottomrule\end{tabular}}
\end{table}
}

\end{document}
